\chapter{Literature Review}

\section{Introduction}

The objective of this literature review is to document and evaluate current knowledge on the use of \acp{LLM}, methods for optimizing the efficiency and scaleability of inference processes, as well as \ac{RAG}-based architectures used with various types of computational environments. With \acp{LLM} continuing to grow in terms of size and complexity, they are resulting in growing computational requirements that necessitate improved methods for deploying and optimizing them locally or distributively (including \ac{HPC}) in order to support increased levels of utilization. At the same time, prior work has demonstrated potential benefits associated with using RAG systems to improve the factual accuracy of responses from \acp{LLM}, reduce instances of hallucination, etc., thus further study will be necessary to investigate the relative performance characteristics of such systems when utilized under a variety of hardware and/or software configurations.

This Chapter brings together prior literature in three broad categories: (i) deployment approaches for \acp{LLM} on Local, Cloud, and \ac{HPC} platforms; (ii) Optimization Approaches to improve Scalability, Memory Usage Reduction and Throughput Improvement of Inference; and (iii) Retrieval Augmentation Systems which include Retrieval Methods, Indexing Techniques, Context Optimizations. In addition, the Chapter provides a review of Benchmarking Standards and Evaluation Methodologies to Assess Computational and Retrieval Performance of \acp{LLM} and \ac{RAG} Pipelines.

Through compilation of findings from these categories, this Chapter identifies several key gaps and shortcomings in the existing body of literature that will guide the experimental design and methodology of this dissertation. Specifically, it draws attention to the lack of studies examining complete \ac{RAG} pipelines within an \ac{HPC} environment, the limited comparative evaluations across different hardware configurations, and the lack of standardization with regard to workflow methodologies for evaluating scalability, retrieval latency, and inference efficiency in \ac{HPC} architectures.

\section{LLM Deployment Approaches}

\subsection{Local and Cloud-Based Deployments}

Research into the deployment of \ac{LLM}s have generally been limited to either cloud based infrastructure, and local systems. The main differences between these types of deployments are related to their available resources (such as CPU/GPU), how scalable they can be (scalability), the total cost of ownership, and the performance characteristics of each environment.

\textbf{Deployments using cloud-based infrastructure} utilize the elastic computing resources provided by major cloud vendors such as Amazon Web Services (\ac{AWS}), Microsoft Azure, and Google Cloud. These services support the use of distributed inference engines, highly scalable \acp{GPU}, and also support the ability to manage model serving applications at scale; thus allowing for reliable execution of very large \ac{LLM} workloads. Additionally studies have shown that through automated scaling, and distributed storage a high degree of efficiency is possible when executing multi-billion parameter models \cite{narayanan2021cloudscale}. Some negative aspects associated with deploying models in the cloud include: the high operating expense (e.g. hourly charges for running the service), Data Transfer Fees (for transferring data across networks), Vendor Lock-In (inability to switch from one cloud provider to another) and Latency (the time it takes for data to travel over the network).

\textbf{Deployments using local infrastructure} require a user to deploy their model(s) to a workstation class \ac{GPU} or on premises server. Local deployments offer several advantages including: predictable performance; low latency per user; and strong data governance. A significant disadvantage of local deployments is the inability to scale up and down to accommodate multiple users or meet the growing demand for increasingly larger \ac{LLM} models.

Given the current trade-offs between control (local deployments) and scalability (cloud deployments), there exists no single deployment platform that currently meets the needs of modern large-scale \ac{LLM} deployments. Therefore, we believe that exploring HPC-based deployments would help to address many of the current disadvantages inherent in both local and cloud-based deployments.

\subsection{HPC-Based Deployments}

High Performance Computing systems offer a potential solution for scaling \ac{LLM} workload due to their highly connected clusters, high-bandwidth interconnects, and massive amounts of total memory. In contrast to cloud based environments optimized for loosely connected service-based applications, \ac{HPC} clusters allow for rapid communication between each node, enabling both efficient distributed training and inference.

A great deal of prior work focuses on distributed training frameworks (Megatron-LM \cite{shoeybi2019megatron}, DeepSpeed's ZeRO \cite{rajbhandari2020deepspeed}) that utilize tensor, pipeline, and data parallelism to enable the training of extremely large models; these approaches also inform inference based on distributed models with memory limitations per device.

In addition to the above references, other researchers have examined aspects of deep learning in \ac{HPC} environments including high-speed interconnects such as InfiniBand and NVLink, the effects of parallel file system \ac{I/O} performance on \ac{LLM}, and the impact of scheduler overhead on performance \cite{abdolrashidi2021large}. However, while many of these studies have identified several of the major factors impacting performance in HPC environments, they typically focused on training rather than inference. Fewer studies have been conducted examining \emph{retrieval}-based workloads which are crucial in \ac{RAG} systems since both latency of the retrieval process and storage throughput affect overall performance of \ac{RAG} systems.

Prior research has shown how various types of system level decisions can be used to optimize different ML workloads using containerization (Apptainer), and scheduling (Slurm). However, few studies have extended this type of evaluation to full \ac{LLM} inference or \ac{RAG} deployments.

While there is clearly an advantage for scaling \ac{LLM} workloads utilizing \ac{HPC} platforms, existing literature identifies significant areas for additional study: limited prior research related to inference performance, minimal prior analysis of vector search and retrieval latency, and no full evaluations of the \ac{RAG} pipeline within \ac{HPC} environments.

\subsection{Model Serving Frameworks}

A number of high-throughput model-serving frameworks have been developed to scale \acp{LLM}:

The \textbf{\ac{vLLM}} framework was the first to use a \textit{PagedAttention} mechanism to improve \ac{KV}-cache efficiency and increase throughput during heavy concurrency as compared to previous mechanisms \cite{kwon2023vllm}. The current version of this framework is likely one of the top performing open source inference engines.

The \textbf{NVIDIA \ac{TIS}} is an inference server framework capable of running multiple frameworks concurrently using dynamic batch sizing and GPU-aware task scheduling. However, it was not initially designed with LLMs in mind, therefore, it has limited ability to optimize LLM-specific parameters.

The \textbf{\ac{TGI})} framework was created for the purpose of optimizing performance when deploying transformer-based text generation models. \ac{TGI} supports tensor-parallelized inference, quantization, and continuous batch processing.

\textbf{RayServe} is a distributed orchestration tool for scaling applications across cluster resources. While RayServe does provide a flexible way to deploy applications across multiple machines, its primary function is focused on the distribution of tasks across nodes rather than providing optimizations for inference speed. 

Therefore, while there are several serving frameworks that can be used to distribute application logic, few have evaluated the impact of hardware constraints found on HPC systems (e.g., batch scheduling, network fabric, and parallel I/O) on the overall performance of these frameworks.

\section{Optimization Techniques for LLM Inference}

\subsection{Quantization and Pruning}

Quantization allows both a reduction in Memory (Memory Usage) and Computation (Computational Complexity), by converting the Weights and Activations from High-Precision (\ac{FP32} / \ac{FP16}) to Lower Precision Formats like \ac{INT8} or \ac{INT4}. Methods including LLM.int8 \cite{dettmers2022int8}, \ac{GPTQ}, and \ac{AWQ} allow for significant reductions in Resource Utilization with minimal losses in Accuracy. This enables Larger Batch Sizes, Faster Inference etc.

Pruning is a method to remove redundant Weights from an existing model to reduce the Size of the model and Computational Complexity. Although early research focused primarily on Vision Models, recent work has applied Pruning to Transformer Models, and has demonstrated that Structured Sparsity can be used to reduce Compute Requirements without large losses in Accuracy.

There are both Advantages and Disadvantages for HPC Environments for Quantizing and Pruning, which will Improve GPU Memory Fragmentation, increase Batching capacity, and decrease I/O Overhead on \acp{PFS}. There are almost no evaluations of these Techniques under the constraints of HPC.

\subsection{Knowledge Distillation}

Distillation has been shown to allow a larger "teacher" LLM to be reduced into a smaller "student". For example, models like DistilBERT \cite{sanh2019distilbert} have demonstrated that many capabilities of an LLM may be preserved in a model that is significantly cheaper than the original model. 

However, the knowledge distillation method will still result in lower latency and memory usage than the teacher model but it does not provide better results on knowledge-intensive tasks. The method also does not inherently prevent hallucination and does not include the ability to ground retrievals in either the data used or the context in which the model was trained. Lastly, the research has not examined how knowledge distillation methods use HPC resources including allocation of those resources, scheduling, and distribution during inference.

\subsection{Parallelization and Sharding}

Strategies to enable \texttt{\ac{DP}}, \texttt{\ac{TP}} and \texttt{\ac{PP}}, represent the most important ways to scale LLMs. \ac{TP} and \ac{PP} as used in Megatron-LM \cite{shoeybi2019megatron} are key to enabling distributed execution. The use of parameter, gradient and optimizer state shards in DeepSpeed ZeRO \cite{rajbhandari2020deepspeed} and PyTorch \ac{FSDP} can be seen as an extension of the ideas behind those two forms of parallelization. 

While there is a substantial body of work analyzing the impact of these methods on the training process, few studies have looked at their impacts when executing inference in a distributed environment, particularly in environments using RAG based pipelines that utilize both retrieval processes and distributed KV-cache implementations along with inter-node communication.

\subsection{Caching and Memory Optimization}

The attention model by FlashAttention \cite { dao2022flashattention } reduces the memory consumption while keeping the accuracy to a high level, thus is very suitable for applications with low memory requirements.

The PagedAttention (vLLM) splits the \ac{KV-cache} to paged blocks, it improves significantly the performance on the systems which serve multiple requests at once.

In addition to sparsification, eviction strategies and cache compressions, researchers are actively exploring other optimizations. However most of them are still in an experimental phase.

Due to multi-user workloads, job schedulers etc., HPC environments have highly varying amounts of available memory. Despite that, there is hardly any literature about applying this type of optimizations on HPC-systems or during multi-node inference.

\section{RAG Systems and Optimization Strategies}

\subsection{Existing Frameworks}

The major frameworks for building RAG pipelines (LangChain \cite{langchain2023}, LlamaIndex \cite{llamaindex2023} and Haystack) offer modules to build the pipeline. Yet according to the literature there are very few benchmarks, they have been tested little on a distributed computing environment and most do not support HPC-specific workflow.

\subsection{Retrieval Efficiency and Hybrid Approaches}

Although, there have been advancements in the area of hybrid retrieval techniques that use a combination of both sparse and dense retrieval strategies to enhance recall and robustness, the research community has made little effort to evaluate retrieval performance with respect to the following categories: 

\begin{itemize}
    \item \acp{PFS},
    \item Multi Node Clusters,
    \item High Performance Computing Interconnects, 
    \item Retrieval performance using Slurm.
\end{itemize}

\subsection{RAG on HPC Environments}

While there has been a relatively small amount of research conducted on RAG in High Performance Computing (HPC) environments, most RAG related research is based upon assumptions that are specific to cloud computing and/or enterprise deployments, where services exist persistently, dynamic autoscaling exists, and container orchestration exists.

In contrast, \ac{HPC} systems utilize batch scheduling (e.g., Slurm) and employ parallel storage systems (e.g. Lustre, \ac{GPFS}). The jobs executed on these systems do not remain active after their completion. These differences create additional complexities regarding how to serve, retrieve, and perform distributed inference.

\section{Benchmarking and Evaluation Studies}

Benchmarks such as \ac{MMLU}, \ac{HELM} \cite{liang2022helm}, and OpenLLM Leaderboard assess \ac{LLM} quality but ignore system performance. \ac{MLPerf} \cite{reddi2020mlperf} measures inference speed but does not evaluate:

\begin{itemize}
    \item autoregressive decoding,
    \item retrieval latency,
    \item distributed inference.
\end{itemize}

RAG-specific benchmarks like \ac{KILT} assess retrieval accuracy but not computational behavior or scalability. Across the board, benchmarks do not capture:

\begin{itemize}
    \item scheduler effects,
    \item inter-node communication,
    \item parallel file system performance,
    \item container overhead.
\end{itemize}

\section{Identified Research Gaps}

The literature reveals several critical gaps:

\begin{itemize}
    \item Lack of systematic studies on \ac{LLM} inference on \ac{HPC} systems.
    \item Lack of research on \ac{RAG} pipelines deployed on \ac{HPC} platforms.
    \item Insufficient analysis of Slurm scheduling and multi-user resource contention.
    \item Limited comparisons of local, cloud, and \ac{HPC} inference performance.
    \item No end-to-end evaluation of vector search and retrieval in \ac{HPC} environments.
    \item Absence of reproducible workflows for LLM/RAG experiments on \ac{HPC}.
\end{itemize}

These gaps prevent practitioners from making informed decisions about infrastructure, optimization choices, and deployment strategies. This thesis addresses these limitations through controlled experiments, cross-environment benchmarking, and analysis of full \ac{RAG} pipeline performance.

\section{Summary}

This chapter reviewed the literature on LLM deployment strategies, inference optimization techniques, RAG System Design, and current evaluation methods for LLM and RAG. While much has been accomplished in scaling and optimizing LLMs, little research is currently being conducted to study deploying and evaluating LLM and RAG workloads on HPC Systems. Also, there exists an unfulfilled need for benchmarking as well as the analysis of HPC-specific issues (e.g., behavior of interconnects, performance of parallel storage, scheduling) that have never before been studied.

Therefore, these shortcomings will be addressed through the methodology presented within this dissertation. Chapter 4 discusses the experimental setup, deployment strategies, optimization techniques, and metrics of evaluation to measure the performance of LLM and RAG across Local, Cloud, and HPC environments.
